\subsection{Neural preference policies and the role of the action oracle}\label{sec:ndu-neural}
We compute the original preference-adjustment economy with the cardinal utility, shock correlation, reflected preference state, and wealth stopping payoff specified above. The main grid has 17 preference points, 25 log-wealth points, and 20 time steps over a unit horizon. Adjustment cost is $k=1$. The expanded control box is
\[
 m=c/X\in[.02,.30],\qquad p\in[0,1.5],\qquad\theta\in[-.45,.45].
\]
The inherited finite-action comparison uses the 175 triples formed by five consumption shares, five portfolio shares, and seven adjustment efforts, including all endpoints. This is the same constrained economy across the matched finite methods, not a change of utility normalization.

Two experiments distinguish action representation from action optimization. First, the R7 finite study trains two-hidden-layer, width-32 tanh critics and categorical actors. At each backward step it computes the complete 175-action payoff vector using the current continuation critic. Its exact argmax provides a label, and the vector also supplies the actor's expected action-loss objective. Thus \emph{the training oracle is exhaustive}, even when the final actor is accurate. It is a useful neural policy-representation experiment, but not evidence that the actor avoided Bellman maximization.

Second, R8 trains a continuous-output actor directly through its own positive-weight payoff. It never receives an all-action value vector or an exact argmax label. The actor receives 160 Adam updates per time level; the continuation critic receives 320 Adam updates and up to 300 L-BFGS iterations, with failed polishing steps rolled back. The actor-free direct comparator optimizes eight independent bounded controls at every state for the same 160 iterations, followed by the same critic fit. Every restart and rejected action query is charged. Neither procedure invokes a Riccati solution, an analytically invariant policy class, or the independently computed reference during training.

\paragraph*{Finite-action neural policies.}
The raw R7 actor is the primary neural representation result. Its full dynamic policy loss is below the declared $.1$ target for each fixed seed. A top-four shortlist and a separately reported correction map further reduce the loss. The correction map is not part of the raw network. The complete finite policy, not only one-step deviations against a fitted critic, is evaluated by independent backward recursion.
\begin{table}[htbp]\centering
\caption{Finite preference policies: R7 evidence}
\begin{small}\begin{tabular}{lrrrrrr}
\toprule
Method & Seed & Raw loss & Final loss & Value error & Seconds & Guard (\%) \\
\midrule
NBO & 11 & 0.0445 & 0.0261 & 0.0924 & 30.59 & 0.294 \\
Direct & 11 & 0.0181 & 0.0181 & 0.0790 & 24.79 & 0.000 \\
NBO & 29 & 0.0500 & 0.0260 & 0.0814 & 30.49 & 0.332 \\
Direct & 29 & 0.0263 & 0.0263 & 0.0551 & 25.19 & 0.000 \\
NBO & 47 & 0.0500 & 0.0263 & 0.1325 & 29.89 & 0.230 \\
Direct & 47 & 0.0264 & 0.0264 & 0.0824 & 25.17 & 0.000 \\
\bottomrule\end{tabular}\end{small}
\par\smallskip\begin{minipage}{.96\linewidth}\footnotesize Loss bounds cover every stored state and time of the same 175-action economy. Seconds include training and the three reported policy verifications; classical backward induction averages 0.102 seconds. A raw actor, shortlist, and corrected policy are distinct objects. Utility losses are absolute units.\end{minipage}
\end{table}

These results repair the earlier R6 actor failures, which remain in the historical numerical record. They do not reverse the cost comparison: the matched direct finite-Bellman fit is faster and more accurate on average, and classical backward induction is substantially cheaper on this small grid. Tables report those comparisons without selecting a favorable seed. The actor's approximation cost is part of the result, not an argument for discarding the economic application.

\paragraph*{Continuous-action optimization without exact labels.}
The initial continuous actor and direct-gradient pilot fail the $.1$ deviation diagnostic. The actor can saturate at a portfolio bound while a lower portfolio share improves continuation value. Four actor heads with output reinitialization do not cure that failure in the retained pilot. We therefore execute a disclosed local-improvement variant: compare the actor proposal with eight box corners, then perform two coordinate-search sweeps at each of eight decreasing step sizes. The same local search is applied to the direct comparator. An actor-free ablation omits all gradient action updates and starts the identical search from feasible proposals. This ablation measures how much of the hybrid result comes from the search itself.
\begin{table}[htbp]\centering
\caption{Continuous-control training without exact maximizing labels}
\begin{small}\begin{tabular}{lrrrrr}
\toprule
Method & Seed & Raw diagnostic & Hybrid gain & Train (s) & Queries ($10^6$) \\
\midrule
Actor + search & 11 & 0.2596 & 0.0617 & 15.01 & 2.261 \\
Actor + search & 29 & 0.0733 & 0.0188 & 28.49 & 2.261 \\
Actor + search & 47 & 0.2597 & 0.0475 & 26.46 & 2.261 \\
Direct + search & 11 & 0.1554 & 0.0303 & 32.02 & 11.841 \\
Direct + search & 29 & 0.1566 & 0.0203 & 32.80 & 11.841 \\
Direct + search & 47 & 0.1554 & 0.0251 & 30.52 & 11.841 \\
Search only & 11 & 5.5370 & 0.0617 & 13.97 & 0.961 \\
Search only & 29 & 7.2093 & 0.0616 & 13.72 & 0.961 \\
Search only & 47 & 8.0316 & 0.0616 & 12.87 & 0.961 \\
\bottomrule\end{tabular}\end{small}
\par\smallskip\begin{minipage}{.96\linewidth}\footnotesize Raw diagnostic is the largest finite-reference payoff minus the raw policy payoff. Hybrid gain is the augmented finite-deviation gain, a lower bound on continuous-action regret, not a certificate. All methods use zero exhaustive training backups and zero exact argmax training labels. Search-only omits action-gradient updates. Training, reference, and full-cover verification costs are separate.\end{minipage}
\end{table}

The table distinguishes three quantities. The raw-network loss against the finite reference tests the actor alone. The augmented finite-deviation gain permits all 175 inherited actions and the learned continuous action at every state and time; it is a \emph{lower bound} on missed continuous-action gains. The complete action-box upper bound in Proposition~\ref{prop:r8-envelope} is a separate verification result. A small augmented gain is not inserted into that proposition as if it were a uniform upper bound. Local changes occur at nearly every active state-time node in the principal hybrid runs, so this experiment does not describe a sparse correction map or an actor-only controller.
\begin{table}[htbp]\centering
\caption{Full continuous-action policy bounds on the two-state nodal economy}
\begin{small}\begin{tabular}{lrrrrr}
\toprule
Policy & Seed & Uniform bound & Center bound & Eval. (s) & Bound $\leq .1$ \\
\midrule
Actor + search & 11 & 0.1294 & 0.0601 & 0.22 & No \\
Actor + search & 29 & 0.1227 & 0.0583 & 0.22 & No \\
Actor + search & 47 & 0.1301 & 0.0608 & 0.25 & No \\
Direct + search & 11 & 0.1153 & 0.0604 & 0.23 & No \\
Direct + search & 29 & 0.1156 & 0.0586 & 0.22 & No \\
Direct + search & 47 & 0.1203 & 0.0609 & 0.25 & No \\
Search only & 11 & 0.1196 & 0.0647 & 0.24 & No \\
Search only & 29 & 0.1179 & 0.0586 & 0.23 & No \\
Search only & 47 & 0.1209 & 0.0610 & 0.22 & No \\
\bottomrule\end{tabular}\end{small}
\par\smallskip\begin{minipage}{.96\linewidth}\footnotesize The full action box is covered at all 425 states and 20 times. The common optimal upper envelope costs 1624.59 seconds across the retained implementations and 40.96 million action-box evaluations. Each row adds its own interval policy evaluation. Unresolved leaves retain their upper bounds. No continuous-state/time error is set to zero.\end{minipage}
\end{table}

The action cover includes binding constraints and wealth exits. Its unresolved boxes remain in the bound when the declared budget is exhausted. The resulting upper array can audit the actor, the direct optimizer, and the local-search ablation without retraining. This produces an operational error account on a multidimensional state grid with a continuous, nonconcave action objective. It does not establish that all runs meet the tighter $.1$ complete-action target, nor that the actor dominates the direct method at matched economic accuracy.

\paragraph*{Alternative representations and economic units.}
Additional same-model comparisons use a gated neural continuation architecture and tensor Chebyshev projection. The gated method uses DGM-style state-dependent gates \citep{sirignano2018}, but here it is fitted to finite Bellman targets, not called a mesh-free differential DGM solver. The projection exploits the same two-state geometry and imposes the same wealth-boundary lifting. Degrees, parameter counts, all-action costs, failures, and actual accuracy appear in the supplement. These are executed post-review comparisons, not results attributed to the earlier unexecuted development protocol. A generic semilinear BSDE comparator is not substituted for this controlled-covariance, stopped problem without reformulating its second-order control term.

All losses are in the specified cardinal utility units. The optimal finite-model value at the center is approximately $-4.8410$, and the initial value range is approximately $[-9.1599,-.1250]$. A loss of $.1$ is about $2.066$ percent of the center's \emph{absolute utility magnitude}; this scaling is not a consumption-equivalent welfare percentage. Endogenous curvature and adjustment costs make a common consumption-equivalent conversion a separate economic calculation.

\paragraph*{Action bounds and controlled refinements.}
The inherited sensitivity study retains both adjustment-cost values, the original 27-action truncation, 175 and 1,053 actions in the expanded box, a wider action box, wider wealth bounds, and finer space--time grids. Consumption and portfolio upper bounds remain frequently active. Expanding the box changes the constrained economy; refining actions inside one box does not. R8 additionally evaluates the \emph{same frozen continuous policy}, bilinearly interpolated in state and held on its original time schedule, on three nested state--time levels. This isolates changes in policy evaluation from retraining. The supplemental differences are observed refinement diagnostics, not an estimated order promoted to a continuous-time error bound. The original continuous control problem, its utility normalization, and its comparative-static questions remain the economic objects of interest.
